\documentclass[11pt,a4paper]{article}
\input{../../sablony/pripravy/priprava.sty}
\usepackage{tikz,pgfplots}
\pgfplotsset{compat=1.18}
\nastavPripravu{1. ročník}{Kvadratické rovnice a nerovnice}{129}{Souhrnná samostatná práce: kvadratické učivo}
\begin{document}
\blok{Souhrnná samostatná práce: kvadratické učivo}
Výběrové procvičování před testem – každá skupina si vybírá úlohy podle chyb a připravenosti.
\blok{Metoda a návaznost na graf}
\begin{itemize}
\item Doporučené minimum: jedna neúplná rovnice, jedna úplná, jedna nerovnice a jedna aplikace.
\item Vyspělejší žáci mohou přidat substituci, absolutní hodnotu a soustavu.
\end{itemize}
\blok{Společně řešené příklady}
\begin{prikladReseniBox}[title={Příklad 1}]
$2x^{2} - 7x + 3=0$
\tcblower
\[
\begin{aligned}
D&=49-4\cdot(2)\cdot(3)=25\\[2pt]
x_1&=\frac{1}{2}\\[2pt]
x_2&=3\\[2pt]
K&=\{\frac{1}{2};3\}
\end{aligned}
\]
\end{prikladReseniBox}
\begin{prikladReseniBox}[title={Příklad 2}]
$x^{2} - 5x + 6\ge0$
\tcblower
\[
\begin{aligned}
f(x)&=x^{2} - 5x + 6\\[2pt]
f(x)=0&\ \Rightarrow\ K_0=\{2;3\}\\[2pt]
K&=(-\infty;2\rangle\cup\langle3;+\infty)
\end{aligned}
\]
\end{prikladReseniBox}
\blok{Samostatné procvičování}
\begin{zadaniReseniBox}[title={Úloha 1}]
$x^{4} - 5x^{2} + 4=0$
\tcblower
\[
\begin{aligned}
t&=x^2\ge0\\[2pt]
0&=t^{2} - 5 \cdot t + 4\\[2pt]
t&\in\{1;4\}\\[2pt]
K&=\{-2;-1;1;2\}
\end{aligned}
\]
\end{zadaniReseniBox}
\begin{zadaniReseniBox}[title={Úloha 2}]
$\left|x^{2} - 4\right|=3$
\tcblower
\[
\begin{aligned}
|f(x)|&=3\\[2pt]
f(x)&=3\quad\vee\quad f(x)=-3\\[2pt]
K_1&=\{- \sqrt{7};\sqrt{7}\}\\[2pt]
K_2&=\{-1;1\}\\[2pt]
K&=\{- \sqrt{7};-1;1;\sqrt{7}\}
\end{aligned}
\]
\end{zadaniReseniBox}
\blok{Volitelná rezerva}
\poznamka{Studenti si vymyslí obdobný příklad se změnou jednoho koeficientu a vysvětlí, jak se změní počet kořenů či řešení.}
\blok{Volitelná výzva -- rovnice s modulem kvadratického výrazu}
\begin{zadaniReseniBox}[title={Rozšiřující příklad}]
$|x^2-5x+6|=2$.
\tcblower
\[\begin{aligned}
x^2-5x+6&=2\ \text{nebo}\ -2\\x^2-5x+4&=0\ \text{nebo}\ x^2-5x+8=0\\D_2&=25-32=-7\\K&=\{1;4\}
\end{aligned}\]
\end{zadaniReseniBox}
\blok{Domácí úkol}
\begin{zadaniReseniBox}[title={Domácí příprava}]
\begin{enumerate}[label=\alph*),itemsep=2pt]
\item $x^{2} - 9=0$
\item $2x^{2} - 7x + 3=0$
\item $x^{2} - 5x + 6<0$
\item $\begin{cases}y=x^{2}\\y=4\end{cases}$
\item $x^{2} - 4 \cdot \left|{x}\right| + 3=0$
\item $\dfrac{\left(x - 3\right) \cdot \left(x + 2\right)}{x - 1}>0$
\end{enumerate}
\tcblower
\begin{enumerate}[label=\alph*),itemsep=2pt]
\item
\[\begin{aligned}
x^2&=9\\|x|&=3\\[2pt]
K&=\{-3;3\}
\end{aligned}\]
\item
\[\begin{aligned}
D&=49-4\cdot(2)\cdot(3)=25\\[2pt]
x_1&=\frac{1}{2}\\[2pt]
x_2&=3\\[2pt]
K&=\{\frac{1}{2};3\}
\end{aligned}\]
\item
\[\begin{aligned}
f(x)&=x^{2} - 5x + 6\\[2pt]
f(x)=0&\ \Rightarrow\ K_0=\{2;3\}\\[2pt]
K&=(2;3)
\end{aligned}\]
\item
\[\begin{aligned}
x^{2}&=4\\[2pt]
0&=x^{2} - 4\\[2pt]
K&=\{[-2;4];\ [2;4]\}
\end{aligned}\]
\item
\[\begin{aligned}
t&=|x|\ge0\\[2pt]
t^{2} - 4 \cdot t + 3&=0\quad\text{(řešíme pro }t\text{)}\\[2pt]
t&\in \{1;3\}\\[2pt]
K&=\{-3;-1;1;3\}
\end{aligned}\]
\item
\[\begin{aligned}
x&\ne 1\quad\text{(podmínka)}\\[2pt]
\text{Kritické body: }&-2,\ 1,\ 3\\[2pt]
K&=(-2;1)\cup(3;+\infty)
\end{aligned}\]
\end{enumerate}
\end{zadaniReseniBox}
\poznamka{Orientační čas: úvod 5 min, společné příklady 18 min, samostatně 15 min, kontrola a shrnutí 7 min. Rezerva je dobrovolná. Číselné parametry modelových aplikací nejsou skutečně naměřená data.}
\end{document}
